\subsection{统计分析协议}
\label{sec:experiment-stats}

所有去噪方法之间的比较评估均采用Wilcoxon符号秩检验，这是一种适用于实验超声波信号指标中观察到的非正态分布的非参数配对检验。检验以双侧配置应用以检测方法对之间的任何系统性差异。

\textbf{多重比较校正。} 当同时比较超过两种方法时，使用Holm--Bonferroni程序~\cite{Holm1979}控制家族误差率。对于$k$个方法比较，$k$个未校正的$p$值排序为$p_{(1)} \leq p_{(2)} \leq \cdots \leq p_{(k)}$；第$i$个排序检验的原假设在$\alpha / (k - i + 1)$水平上被拒绝。这一逐步降级程序比标准Bonferroni校正具有更高的统计功效。

\textbf{显著性阈值。} 所有检验采用$p < 0.05$的显著性水平。

\textbf{报告规范。} 指标值报告为跨$N$个试件的median $\pm$四分位距(IQR)，或当分布近似对称且样本量足以评估正态性时（Shapiro--Wilk检验，$p > 0.05$）报告为均值$\pm$标准差。每个试件贡献$M$个测量位置；指标首先在试件内平均，然后计算跨试件的汇总统计。

\textbf{效应量。} 在识别出统计显著差异时，使用等级双列相关系数$r = Z / \sqrt{N}$报告效应量，其中$Z$是Wilcoxon符号秩检验的标准化检验统计量，$N$是配对观察的数量。效应量按照Cohen的惯例解释为小（$|r| = 0.1$）、中（$|r| = 0.3$）或大（$|r| = 0.5$）。